% Spectre infrarouge d'un liquide de formule brute C2H6O (exercice 5, chapitre 11)
\begin{tikzpicture}
\begin{axis}[otIR,width=104mm,height=44mm]
\addplot[otSpectre] {100 \bandeIR{3350}{170}{62} \bandeIR{2950}{55}{48} \bandeIR{1420}{55}{22} \bandeIR{1080}{45}{55} \bandeIR{880}{30}{18}};
\draw[pcRepere] (axis cs:3350,40) -- (axis cs:3350,76);
\node[pcLegendeG,fill=white,inner sep=1.5pt,draw=otNavy,circle,text=otNavy] at (axis cs:3350,83) {1};
\draw[pcRepere] (axis cs:2950,54) -- (axis cs:2600,80);
\node[pcLegendeG,fill=white,inner sep=1.5pt,draw=otNavy,circle,text=otNavy] at (axis cs:2520,86) {2};
\end{axis}
\end{tikzpicture}
\par\smallskip\legende{Le spectre infrarouge du liquide étudié ; les bandes 1 et 2 sont à attribuer.}
